\documentclass[12pt]{article}
\usepackage{graphicx}
\usepackage{geometry}
\usepackage{listings}
\usepackage[T1]{fontenc}
\usepackage{url}
\usepackage[hidelinks]{hyperref}
\usepackage{float}
\usepackage{caption}
\usepackage{bookmark}
\geometry{margin=1in}

\lstset{
  basicstyle=\ttfamily\small,
  breaklines=true,
  breakatwhitespace=true,
  columns=fullflexible
}

% ---- Macro for bottom-only cropping ----
% Usage: \cropbottom[<bottom>]{<file>}
% Example: \cropbottom[150pt]{step1.png}
\newcommand{\cropbottom}[2][150pt]{%
  \includegraphics[width=0.9\textwidth, trim={0 #1 0 0}, clip]{#2}%
}

\begin{document}

\begin{center}
    {\Large \textbf{COMP 324 Network Security}} \\[4pt]
    {\large Module 3 Lab 2: Permission Manipulation} \\[10pt]
    \textbf{Lab Report} \\[6pt]
    \textbf{Student: Erik Gonzalez} \\[4pt]
    \textbf{Professor: Samuel Addington} \\[4pt]
    \textbf{Date: \today}
\end{center}

\section*{Step 1: Changing File Permissions}
\begin{figure}[H]
    \centering
    \cropbottom[800pt]{step1.png}
    
\end{figure}

\noindent \textbf{Explanation:}  
The file was first returned to \texttt{student1} ownership using \texttt{chown}.  
\texttt{chmod 444} set the file to read-only for all users,  
append failed, verifying read-only enforcement.  
Next, \texttt{chmod 600} gave read/write access only to the owner.  
Verification with \texttt{ls -la} showed \texttt{-rw-------}, and appending content succeeded, confirming the permission change.

\section*{Step 2: Special Permissions}
\begin{figure}[H]
    \centering
    \cropbottom[700pt]{step2.png}
    
\end{figure}
\noindent \textbf{Explanation:}  


\texttt{chmod +x} added execute permission.  
\texttt{chmod 4755} set SUID, changing owner execute to \texttt{s}.  
\texttt{chmod 2755} set SGID, changing group execute to \texttt{s}.  
\texttt{chmod 1755} set the Sticky Bit, adding \texttt{t} in the others’ field.  
\noindent \textbf{Observations:}  
\begin{itemize}
    \item The SUID bit replaces the owner’s execute bit with \texttt{s}, making programs run with file owner privileges.
    \item SGID on a directory ensures all new files inherit the directory’s group    \item The Sticky Bit on a directory prevents users from deleting others’ files\end{itemize}

\section*{Step 3: Access Control Lists (ACLs)}
\begin{figure}[H]
    \centering
    \cropbottom{step3.png}
    \end{figure}

\noindent \textbf{Explanation:}  
  
A user ACL entry was added with \texttt{setfacl -m u:student1:r}, with read access, and a group ACL entry with \texttt{setfacl -m g:developers:w}, granting write access.  
\texttt{getfacl} showed the additional ACL entries alongside the standard owner/group/other fields.  
ACLs were removed with \texttt{setfacl -b}, restoring default permissions.

\noindent \textbf{Output:}  
\begin{lstlisting}
# file: home/student1/project/test.txt
# owner: student1
# group: student1
user::rw-
user:student1:r--
group::---
group:developers:w-
mask::rw-
other::---
\end{lstlisting}
This confirms that ACLs can add specific user and group permissions independent of the standard permission bits.

\section*{Conclusion}
This lab showed me how changing permissions and ownership directly controls what users can do with files.  
I learned how bits like SUID, SGID, and the Sticky Bit change behavior, and how ACLs give more fine-grained control.  
Overall, it made permissions feel a lot clearer and less confusing.  \end{document}

